\documentclass[UTF8,12pt]{ctexart}
\usepackage[a4paper,margin=24mm]{geometry}
\usepackage{amsmath,amssymb,bm,graphicx,adjustbox}
\newcommand{\fitmath}[1]{\begin{adjustbox}{max width=\linewidth}$\displaystyle #1$\end{adjustbox}}
\setlength{\parindent}{0pt}
\setlength{\parskip}{.7em}
\sloppy
\allowdisplaybreaks
\begin{document}
\section*{2004 年考研数学三 · 真题}
\subsection*{2004年全国硕士研究生招生考试试题}

\subsection*{一、填空题（本题共6小题，每小题4分，满分24分）}

（1）若 \(\displaystyle \lim_{x\to0}\dfrac{\displaystyle \sin x}{\displaystyle e^x-a}(\cos x-b)=5\)，则 \(\displaystyle a=\)\_\_\_\_\_\_，\(\displaystyle b=\)\_\_\_\_\_\_。


（2）函数 \(\displaystyle f(u,\allowbreak v)\) 由关系式 \(\displaystyle f[xg(y),\allowbreak y]=x+g(y)\) 确定，其中函数 \(\displaystyle g(y)\) 可微，且 \(\displaystyle g(y)\ne0\)，则 \(\displaystyle \dfrac{\displaystyle \partial^2f}{\displaystyle \partial u\partial v}=\)\_\_\_\_\_\_。


（3）设 \(\displaystyle f(x)=\begin{cases}xe^{x^2},\allowbreak &-\dfrac{\displaystyle 1}{\displaystyle 2}\le x<\dfrac{\displaystyle 1}{\displaystyle 2},\allowbreak \\-1,\allowbreak &x\ge\dfrac{\displaystyle 1}{\displaystyle 2},\allowbreak \end{cases}\) 则 \(\displaystyle \int_{1/2}^{1}f(x-1)\,\mathrm dx=\)\_\_\_\_\_\_。


（4）二次型 \(\displaystyle f(x_1,\allowbreak x_2,\allowbreak x_3)=(x_1+x_2)^2+(x_2-x_3)^2+(x_3+x_1)^2\) 的秩为\_\_\_\_\_\_。


（5）设随机变量 \(\displaystyle X\) 服从参数为 \(\displaystyle \lambda\) 的指数分布，则 \(\displaystyle P\{X>\sqrt{D(X)}\}=\)\_\_\_\_\_\_。


（6）设总体 \(\displaystyle X\) 服从正态分布 \(\displaystyle N(\mu_1,\allowbreak \sigma^2)\)，总体 \(\displaystyle Y\) 服从正态分布 \(\displaystyle N(\mu_2,\allowbreak \sigma^2)\)，\(\displaystyle X_1,\allowbreak \ldots,\allowbreak X_{n_1}\) 和 \(\displaystyle Y_1,\allowbreak \ldots,\allowbreak Y_{n_2}\) 分别是来自总体 \(\displaystyle X\) 和 \(\displaystyle Y\) 的简单随机样本，则 \(\displaystyle E\left[\dfrac{\displaystyle \sum_{i=1}^{n_1}(X_i-\bar X)^2+\displaystyle \sum_{j=1}^{n_2}(Y_j-\bar Y)^2}{\displaystyle n_1+n_2-2}\right]=\)\_\_\_\_\_\_。


\subsection*{二、选择题（本题共8小题，每小题4分，满分32分）}

（7）函数 \(\displaystyle f(x)=\dfrac{\displaystyle |x|\sin(x-2)}{\displaystyle x(x-1)(x-2)^2}\) 在下列哪个区间内有界（　）。
（A）\(\displaystyle (-1,\allowbreak 0)\)。（B）\(\displaystyle (0,\allowbreak 1)\)。（C）\(\displaystyle (1,\allowbreak 2)\)。（D）\(\displaystyle (2,\allowbreak 3)\)。


（8）设 \(\displaystyle f(x)\) 在 \(\displaystyle (-\infty,\allowbreak +\infty)\) 内有定义，且 \(\displaystyle \lim_{x\to\infty}f(x)=a\)，\(\displaystyle g(x)=\begin{cases}f(1/x),\allowbreak &x\ne0,\allowbreak \\0,\allowbreak &x=0,\allowbreak \end{cases}\) 则（　）。
（A）\(\displaystyle x=0\) 必是 \(\displaystyle g(x)\) 的第一类间断点。（B）\(\displaystyle x=0\) 必是 \(\displaystyle g(x)\) 的第二类间断点。（C）\(\displaystyle x=0\) 必是 \(\displaystyle g(x)\) 的连续点。（D）\(\displaystyle g(x)\) 在点 \(\displaystyle x=0\) 处的连续性与 \(\displaystyle a\) 的取值有关。


（9）设 \(\displaystyle f(x)=|x(1-x)|\)，则（　）。
（A）\(\displaystyle x=0\) 是 \(\displaystyle f(x)\) 的极值点，但 \(\displaystyle (0,\allowbreak 0)\) 不是曲线 \(\displaystyle y=f(x)\) 的拐点。（B）\(\displaystyle x=0\) 不是 \(\displaystyle f(x)\) 的极值点，但 \(\displaystyle (0,\allowbreak 0)\) 是曲线 \(\displaystyle y=f(x)\) 的拐点。（C）\(\displaystyle x=0\) 是 \(\displaystyle f(x)\) 的极值点，且 \(\displaystyle (0,\allowbreak 0)\) 是曲线 \(\displaystyle y=f(x)\) 的拐点。（D）\(\displaystyle x=0\) 不是 \(\displaystyle f(x)\) 的极值点，\(\displaystyle (0,\allowbreak 0)\) 也不是曲线 \(\displaystyle y=f(x)\) 的拐点。


（10）设有以下命题：\textcircled{1} 若 \(\displaystyle \sum_{n=1}^{\infty}(u_{2n-1}+u_{2n})\) 收敛，则 \(\displaystyle \sum_{n=1}^{\infty}u_n\) 收敛；\textcircled{2} 若 \(\displaystyle \sum_{n=1}^{\infty}u_n\) 收敛，则 \(\displaystyle \sum_{n=1}^{\infty}u_{n+100}\) 收敛；


（10）（续）\textcircled{3} 若 \(\displaystyle \lim_{n\to\infty}\dfrac{\displaystyle u_{n+1}}{\displaystyle u_n}>1\)，则 \(\displaystyle \sum_{n=1}^{\infty}u_n\) 发散；\textcircled{4} 若 \(\displaystyle \sum_{n=1}^{\infty}(u_n+v_n)\) 收敛，则 \(\displaystyle \sum_{n=1}^{\infty}u_n\)、\(\displaystyle \sum_{n=1}^{\infty}v_n\) 都收敛。则以上命题中正确的是（　）。（A）\textcircled{1}\textcircled{2}。（B）\textcircled{2}\textcircled{3}。（C）\textcircled{3}\textcircled{4}。（D）\textcircled{1}\textcircled{4}。


（11）设 \(\displaystyle f^{\prime}(x)\) 在 \(\displaystyle [a,\allowbreak b]\) 上连续，且 \(\displaystyle f^{\prime}(a)>0\)，\(\displaystyle f^{\prime}(b)<0\)，则下列结论中错误的是（　）。
（A）至少存在一点 \(\displaystyle x_0\in(a,\allowbreak b)\)，使得 \(\displaystyle f(x_0)>f(a)\)。（B）至少存在一点 \(\displaystyle x_0\in(a,\allowbreak b)\)，使得 \(\displaystyle f(x_0)>f(b)\)。（C）至少存在一点 \(\displaystyle x_0\in(a,\allowbreak b)\)，使得 \(\displaystyle f^{\prime}(x_0)=0\)。（D）至少存在一点 \(\displaystyle x_0\in(a,\allowbreak b)\)，使得 \(\displaystyle f(x_0)=0\)。


（12）设 \(\displaystyle n\) 阶矩阵 \(\displaystyle A\) 与 \(\displaystyle B\) 等价，则必有（　）。
（A）当 \(\displaystyle |A|=a\;(a\ne0)\) 时，\(\displaystyle |B|=a\)。（B）当 \(\displaystyle |A|=a\;(a\ne0)\) 时，\(\displaystyle |B|=-a\)。（C）当 \(\displaystyle |A|\ne0\) 时，\(\displaystyle |B|=0\)。（D）当 \(\displaystyle |A|=0\) 时，\(\displaystyle |B|=0\)。


（13）设 \(\displaystyle n\) 阶矩阵 \(\displaystyle A\) 的伴随矩阵 \(\displaystyle A^*\ne O\)，若 \(\displaystyle \xi_1,\allowbreak \xi_2,\allowbreak \xi_3,\allowbreak \xi_4\) 是非齐次线性方程组 \(\displaystyle Ax=b\) 的互不相等的解，则对应的齐次线性方程组 \(\displaystyle Ax=0\) 的基础解系（　）。
（A）不存在。（B）仅含一个非零解向量。（C）含有两个线性无关的解向量。（D）含有三个线性无关的解向量。


（14）设随机变量 \(\displaystyle X\) 服从正态分布 \(\displaystyle N(0,\allowbreak 1)\)，对给定的 \(\displaystyle \alpha\;(0<\alpha<1)\)，数 \(\displaystyle u_\alpha\) 满足 \(\displaystyle P\{X>u_\alpha\}=\alpha\)。若 \(\displaystyle P\{|X|<x\}=\alpha\)，则 \(\displaystyle x\) 等于（　）。
（A）\(\displaystyle u_{\alpha/2}\)。（B）\(\displaystyle u_{1-\alpha/2}\)。（C）\(\displaystyle u_{(1-\alpha)/2}\)。（D）\(\displaystyle u_{1-\alpha}\)。


\subsection*{三、解答题（本题共9小题，满分94分。解答应写出文字说明、证明过程或演算步骤）}

（15）（本题满分8分）求 \(\displaystyle \lim_{x\to0}\left(\dfrac{\displaystyle 1}{\displaystyle \sin^2x}-\dfrac{\displaystyle \cos^2x}{\displaystyle x^2}\right)\)。


（16）（本题满分8分）求 \(\displaystyle \iint_D(\sqrt{x^2+y^2}+y)\,\mathrm d\sigma\)，其中 \(\displaystyle D\) 是由圆 \(\displaystyle x^2+y^2=4\) 和 \(\displaystyle (x+1)^2+y^2=1\) 所围成的平面区域。


（17）（本题满分8分）设 \(\displaystyle f(x),\allowbreak g(x)\) 在 \(\displaystyle [a,\allowbreak b]\) 上连续，且满足 \(\displaystyle \int_a^xf(t)\,\mathrm dt\ge\displaystyle \int_a^xg(t)\,\mathrm dt\; (x\in[a,\allowbreak b))\)，\(\displaystyle \int_a^bf(t)\,\mathrm dt=\displaystyle \int_a^bg(t)\,\mathrm dt\)，证明 \(\displaystyle \int_a^bxf(x)\,\mathrm dx\le\displaystyle \int_a^bxg(x)\,\mathrm dx\)。


（18）（本题满分9分）设某商品的需求函数为 \(\displaystyle Q=100-5P\)，其中价格 \(\displaystyle P\in(0,\allowbreak 20)\)，\(\displaystyle Q\) 为需求量。（I）求需求量对价格的弹性 \(\displaystyle E_d\;(E_d>0)\)；（II）推导 \(\displaystyle \dfrac{\displaystyle \mathrm dR}{\displaystyle \mathrm dP}=Q(1-E_d)\)（其中 \(\displaystyle R\) 为收益），并用弹性 \(\displaystyle E_d\) 说明价格在何范围内变化时，降低价格反而使收益增加。


（19）（本题满分9分）设级数 \(\displaystyle \dfrac{\displaystyle x^4}{\displaystyle 2\times4}+\dfrac{\displaystyle x^6}{\displaystyle 2\times4\times6}+\dfrac{\displaystyle x^8}{\displaystyle 2\times4\times6\times8}+\cdots\;(-\infty<x<+\infty)\) 的和函数为 \(\displaystyle S(x)\)。求：（I）\(\displaystyle S(x)\) 所满足的一阶微分方程；（II）\(\displaystyle S(x)\) 的表达式。


（20）（本题满分13分）设 \(\displaystyle \alpha_1=(1,\allowbreak 2,\allowbreak 0)^{\mathrm T}\)，\(\displaystyle \alpha_2=(1,\allowbreak a+2,\allowbreak -3a)^{\mathrm T}\)，\(\displaystyle \alpha_3=(-1,\allowbreak -b-2,\allowbreak a+2b)^{\mathrm T}\)，\(\displaystyle \beta=(1,\allowbreak 3,\allowbreak -3)^{\mathrm T}\)，试讨论当 \(\displaystyle a,\allowbreak b\) 为何值时，（I）\(\displaystyle \beta\) 不能由 \(\displaystyle \alpha_1,\allowbreak \alpha_2,\allowbreak \alpha_3\) 线性表示；（II）\(\displaystyle \beta\) 可由它们唯一地线性表示，并求出表示式；（III）\(\displaystyle \beta\) 可由它们线性表示，但表示式不唯一，并求出表示式。


（21）（本题满分13分）设 \(\displaystyle n\) 阶矩阵 \(\displaystyle A=\begin{pmatrix}1&b&\cdots&b\\b&1&\cdots&b\\\vdots&\vdots&&\vdots\\b&b&\cdots&1\end{pmatrix}\)。（I）求 \(\displaystyle A\) 的特征值和特征向量；（II）求可逆矩阵 \(\displaystyle P\)，使得 \(\displaystyle P^{-1}AP\) 为对角矩阵。


（22）（本题满分13分）设 \(\displaystyle A,\allowbreak B\) 为两个随机事件，且 \(\displaystyle P(A)=\dfrac{\displaystyle 1}{\displaystyle 4}\)，\(\displaystyle P(B\mid A)=\dfrac{\displaystyle 1}{\displaystyle 3}\)，\(\displaystyle P(A\mid B)=\dfrac{\displaystyle 1}{\displaystyle 2}\)。令 \(\displaystyle X=\begin{cases}1,\allowbreak &A\text{发生},\allowbreak \\0,\allowbreak &A\text{不发生},\allowbreak \end{cases}\quad Y=\begin{cases}1,\allowbreak &B\text{发生},\allowbreak \\0,\allowbreak &B\text{不发生}.\end{cases}\) 求：（I）二维随机变量 \(\displaystyle (X,\allowbreak Y)\) 的概率分布；（II）\(\displaystyle X\) 与 \(\displaystyle Y\) 的相关系数 \(\displaystyle \rho_{XY}\)；（III）\(\displaystyle Z=X^2+Y^2\) 的概率分布。


（23）（本题满分13分）设随机变量 \(\displaystyle X\) 的分布函数为 \(\displaystyle F(x;\alpha,\allowbreak \beta)=\begin{cases}1-(\dfrac{\displaystyle \alpha}{\displaystyle x})^\beta,\allowbreak &x>\alpha,\allowbreak \\0,\allowbreak &x\le\alpha,\allowbreak \end{cases}\) 其中参数 \(\displaystyle \alpha>0\)，\(\displaystyle \beta>1\)。设 \(\displaystyle X_1,\allowbreak X_2,\allowbreak \ldots,\allowbreak X_n\) 为来自总体 \(\displaystyle X\) 的简单随机样本。（I）当 \(\displaystyle \alpha=1\) 时，求未知参数 \(\displaystyle \beta\) 的矩估计量；（II）当 \(\displaystyle \alpha=1\) 时，求未知参数 \(\displaystyle \beta\) 的最大似然估计量；（III）当 \(\displaystyle \beta=2\) 时，求未知参数 \(\displaystyle \alpha\) 的最大似然估计量。

\end{document}
